\documentclass[aspectratio=169]{beamer}
\usepackage[no-math,deluxe,haranoaji]{luatexja-preset}
\renewcommand{\kanjifamilydefault}{\gtdefault}
\renewcommand{\emph}[1]{{\upshape\bfseries #1}}
\usetheme{metropolis}
\metroset{block=fill}
\setbeamertemplate{navigation symbols}{}
\setbeamertemplate{blocks}[rounded][shadow=false]
\usecolortheme[rgb={0.7,0.2,0.2}]{structure}
%%%%%%%%%%%%%%%%%%%%%%%%%%%
\usepackage{media9}
%%%%%%%%%%%%%%%%%%%%%%%%%%%
%% さまざまなアイコン
%%%%%%%%%%%%%%%%%%%%%%%%%%%
\usepackage{fontawesome}
\usepackage{figchild}
\usepackage{twemojis}
\usepackage{utfsym}
\usepackage{bclogo}
\usepackage{marvosym}
\usepackage{fontmfizz}
\usepackage{tipa}
\usepackage{circledsteps}
%%%%%%%%%%%%%%%%%%%%%%%%%%%
\usepackage{tikz}
\usetikzlibrary{backgrounds}
\usepackage{tcolorbox}
\usepackage{myascolorbox} %ascolorbox.sty(https://github.com/yasunari/ascolorbox)をlualatexで使えるようにzw2箇所を\zwに改変したもの
\usetikzlibrary{tikzmark}
\usepackage{tikzpeople}
\usepackage[dvipsnames]{xcolor}
\usepackage{amsmath}
\usepackage{tabularray}
\UseTblrLibrary{booktabs}
%%%%%%%%%%%%%%%%%%%%%%%%%%%
%% 場合分け
\usepackage{cases}
%%%%%%%%%%%%%%%%%%%%%%%%%%%
% \myAnch{<名前>}{<色>}{<テキスト>}
% 指定のテキストを指定の色の四角枠で囲み, 指定の名前をもつTikZの
% ノードとして出力する. 図には remeber picture 属性を付けている
% ので外部から参照可能である.
\newcommand*{\myAnch}[3]{%
  \tikz[remember picture,baseline=(#1.base)]
    \node[draw,rectangle,#2] (#1) {\normalcolor #3};
}
%%%%%%%%%%%%%%%%%%%%%%%%%%%%
%% 音声リンク表示
\newcommand{\myaudio}[1]{\href{#1}{\faVolumeUp}}
%%%%%%%%%%%%%%%%%%%%%%%%%%%
% \myEmph コマンドの定義
%\newcommand{\myEmph}[3]{%
%    \textbf<#1>{\color<#1>{#2}{#3}}%
%}
\usepackage{xparse} % xparseパッケージの読み込み
\NewDocumentCommand{\myEmph}{O{} m m}{%
    \def\argOne{#1}%
    \ifx\argOne\empty
        \textbf{\color{#2}{#3}}% オプション引数が省略された場合
    \else
        \textbf<#1>{\color<#1>{#2}{#3}}% オプション引数が指定された場合
    \fi
}
%%%%%%%%%%%%%%%%%%%%%%%%%%%
%% 文末の上昇イントネーション記号 \myRisingPitch
%% 通常のイントネーション \myDownwardPitch
%% https://note.com/dan_oyama/n/n8be58e8797b2
%%%%%%%%%%%%%%%%%%%%%%%%%%%
\newcommand{\myRisingPitch}{
\begin{tikzpicture}[scale=0.3,baseline=0.3]
\draw[->,>=stealth] (0,0) to[bend right=45] (1,1);
\end{tikzpicture}
}
\newcommand{\myDownwardPitch}{
\begin{tikzpicture}[scale=0.3,baseline=0.3]
\draw[->,>=stealth] (0,1) to[bend left=45] (1,0);
\end{tikzpicture}
}
%%%%%%%%%%%%%%%%%%%%%%%%%%%
\title{English is fun.}
\subtitle{Are you hungry? --- Yes, I am. / No, I'm not.}
\author{}
\institute[]{}
\date[]

%%%%%%%%%%%%%%%%%%%%%%%%%%%%
%% TEXT
%%%%%%%%%%%%%%%%%%%%%%%%%%%%
\begin{document}
\begin{frame}[plain]
  \titlepage
\end{frame}

\section*{授業の流れ}
\begin{frame}[plain]
  \frametitle{授業の流れ}
  \tableofcontents
\end{frame}

%%%%%%%%%%%%%%%%%%%%%%%%%%%%
\section{復習}
\subsection{be動詞の疑問文}
%%%%%%%%%%%%%%%%%%%%%%%%%%%
\begin{frame}[plain]{要点}

be動詞を用いた文の疑問文のつくり方
\begin{block}<2->{Topics for Today}
\begin{itemize}\setbeamertemplate{items}[square]
 \item<3->   be動詞を文の先頭に置く\\
\hfill{\scriptsize He is rich. $\rightarrow$ Is he rich?}
 \item<4->   文末に`?'をつける\\
\hfill{}\visible<5->{{ \scriptsize `?'をquestion mark（クエスチョンマーク）といいます}}
 \item<6->   イントネーションは\myRisingPitch
\end{itemize}
     \end{block}
\end{frame}
%%%%%%%%%%%%%%%%%%%%%%%%%%%%%%
\section{be動詞の疑問文への答え方}
\subsection{Are you 〜 ? と聞かれたら}
%%%%%%%%%%%%%%%%%%%%%%%%%%%%%%
 \begin{frame}[plain,t]{Are you 〜 ? と聞かれたら}
 \large

\visible<2->{Are you from New York?}

\vspace{20pt}

\mbox{}\hspace{100pt}%
\visible<3->{$\left\{\begin{tabular}{l}
         \visible<3->{\text{Yes, I am.}}\\
         \visible<4->{\alt<5,6,7>{\text{No, \Circled{I am} not.}}{\text{No, I am not.}}}\\
         \visible<5->{\text{(}= \text{No, I'm not.)}}
        \end{tabular}\right.$}

\hfill\visible<6->{{\small NoのときI amを縮めてNo, \textcolor{Maroon}{\bfseries I'm} not.ともいいます（短縮形）}}\\
\hfill\visible<7->{{\small でもNo, I \textcolor{Maroon}{\bfseries amn't}.とはいいません}}

\hfill{\tiny 0117}\,{\scriptsize \myaudio{./audio/009_answer_be_01.mp3}}
\end{frame}
%%%%%%%%%%%%%%%%%%%%%%%%%%%%%%%%
\subsection{Is he 〜 ?  / Is she 〜 ? と聞かれたら}
%%%%%%%%%%%%%%%%%%%%%%%%%%%%%%%
\begin{frame}[plain,t]{Is he  〜 ? / Is she 〜 ? と聞かれたら}
 \large

\begin{columns}[t]
\begin{column}{.475\textwidth}
\visible<2->{Is he a baseball player?}

\vspace{20pt}

\mbox{}\hspace{40pt}%
\visible<3->{$\left\{\begin{tabular}{l}
         \visible<3->{\text{Yes, he is.}}\\
         \visible<4->{\temporal<5>{\text{No, he is not.}}{\text{No, \Circled{he is} not.}}{\text{No, he \Circled{is not.}}}}\\
         \visible<5->{\text{(}= \text{No, he's not.)}}\\
         \visible<6->{\text{(}= \text{No, he isn't.)}}
       \end{tabular}\right.$}

\vspace{10pt}

\hfill\visible<5->{{\small Noのときhe isを縮めてNo, \textcolor{Maroon}{\bfseries he's} not.}}

\vspace{-5pt}

\hfill\visible<6->{{\small \phantom{Noのとき}is notを縮めてNo, he \textcolor{Maroon}{\bfseries isn't.}}}

\end{column}
\begin{column}{.475\textwidth}
\visible<7->{Is she a good singer?}

\vspace{20pt}

\mbox{}\hspace{40pt}
\visible<8->{$\left\{\begin{tabular}{l}
         \visible<8->{\text{Yes, she is.}}\\
         \visible<9->{\temporal<10>{\text{No, she is not.}}{\text{No, \Circled{she is} not.}}{\text{No, she \Circled{is not.}}}}\\
         \visible<10->{\text{(}= \text{No, she's not.)}}\\
         \visible<11->{\text{(}= \text{No, she isn't.)}}
       \end{tabular}\right.$}


\vspace{20pt}

\hfill{\tiny 0341}\,{\scriptsize \myaudio{./audio/009_answer_be_02.mp3}}

\end{column}
\end{columns}
\end{frame}
%%%%%%%%%%%%%%%%%%%%%%%%%%%%
\subsection{Is John 〜 ?  / Is Emily 〜 ? と聞かれたら}
%%%%%%%%%%%%%%%%%%%%%%%%%%%
\begin{frame}[plain,t]{Is John  〜 ? / Is Emily 〜 ? と聞かれたら}
 \large

\begin{columns}[t]
\begin{column}{.49\textwidth}
\visible<2->{Is \myAnch{john}{Maroon}{John} a teacher?}

\vspace{20pt}

\mbox{}\hfill%
\visible<3->{$\left\{\begin{tabular}{l}
         \visible<3->{Yes, \myAnch{he1}{Maroon}{he} is.}\\
         \visible<4->{No, \myAnch{he2}{Maroon}{he} is not.}\\
         \visible<5->{\text{(}= \text{No, he's not.)}}\\
         \visible<6->{\text{(}= \text{No, he isn't.)}}
        \end{tabular}\right.$}

\begin{tikzpicture}[remember picture,overlay]
 \visible<3->{\draw[->,line width=2pt,opacity=.5,Maroon] (john.south) to[out=-30,in=160] (he1.north west);}
 \visible<4->{\draw[->,line width=2pt,opacity=.5,Maroon] (john.south) to[out=-110, in=155] ([yshift=3pt]he2.west);}
\end{tikzpicture}

\end{column}
%%%%%%%%%%%%%%%%%%%%
\begin{column}{.49\textwidth}
\visible<7->{Is \myAnch{emily}{Maroon}{Emily} from Sydney?\\
\hfill{\scriptsize Sydney \textipa{/s\'Idni/}}}

\vspace{20pt}

\mbox{}\hfill%
\visible<8->{$\left\{\begin{tabular}{l}
         \visible<8->{Yes, \myAnch{she1}{Maroon}{she} is.}\\
         \visible<9->{No, \myAnch{she2}{Maroon}{she} is not.}\\
         \visible<10->{\text{(}= \text{No, she's not.)}}\\
         \visible<11->{\text{(}= \text{No, she isn't.)}}
        \end{tabular}\right.$}

{\begin{tikzpicture}[remember picture,overlay]
 \visible<8->{\draw[->,line width=2pt,opacity=.5,Maroon] (emily.south) to[out=-30,in=160] (she1.north west);}
 \visible<9->{\draw[->,line width=2pt,opacity=.5,Maroon] (emily.south) to[out=-165, in=165] ([yshift=3pt]she2.west);}
\end{tikzpicture}}

\vspace{20pt}

\hfill{\tiny 0340}\,{\scriptsize \myaudio{./audio/009_answer_be_03.mp3}}
\end{column}
\end{columns}

\visible<3->{\small 答えるときは、名前そのものではなくheやsheを用います}
\end{frame}
%%%%%%%%%%%%%%%%%%%%%%%%%%%%%%%%%%%%%%%%%%%
 \subsection{Are they 〜 ? と聞かれたら}
%%%%%%%%%%%%%%%%%%%%%%%%%%%%%%%%%%%%%%%%%%
\begin{frame}[plain,t]{Are they 〜 ? と聞かれたら}
 \large

\visible<2->{Are they your friends?\hfill{\scriptsize friend \textipa{/fr\'end/}}}

\vspace{10pt}

\mbox{}\hspace{100pt}%
\visible<3->{$\left\{\begin{tabular}{l}
         \visible<3->{Yes, they are.}\\
         \visible<4->{\temporal<5>{\text{No, they are not.}}{\text{No, \Circled{they are} not.}}{\text{No, they \Circled{are not.}}}}\\
         \visible<5->{\text{(}= \text{No, they're not.)}}\\
         \visible<6->{\text{(}= \text{No, they aren't.)}}
        \end{tabular}\right.$}

\vspace{10pt}

\hfill\visible<5->{{\small Noのときthey areを縮めてNo, \textcolor{Maroon}{they're} not.}}

\vspace{-8pt}

\hfill\visible<6->{{\small \phantom{Noのとき}are notを縮めてNo, they \textcolor{Maroon}{aren't.}}}

\vspace{20pt}

\hfill{\tiny 0155}\,{\scriptsize \myaudio{./audio/009_answer_be_04.mp3}}
\end{frame}
%%%%%%%%%%%%%%%%%%%%%%%%%%%%%%%%
\begin{frame}[plain,t]{Are A and B 〜 ? と聞かれたら}
 \large

\visible<2->{Are \myAnch{s1}{Maroon}{George and Peter} your friends?}

\vspace{10pt}

\mbox{}\hspace{100pt}%
\visible<3->{$\left\{\begin{tabular}{l}
         \visible<3->{Yes, \myAnch{they4}{Maroon}{they} are.}\\
         \visible<4->{No, \myAnch{they5}{Maroon}{they} are not.}\\
         \visible<5->{\text{(}= \text{No, they're not.)}}\\
         \visible<6->{\text{(}= \text{No, they aren't.)}}
       \end{tabular}\right.$}

\begin{tikzpicture}[remember picture,overlay]
 \visible<3->{\draw[->,line width=2pt,opacity=.5,Maroon] (s1.south) to[out=-30,in=160] (they4.north west);}
 \visible<4->{\draw[->,line width=2pt,opacity=.5,Maroon] (s1.south) to[out=-90, in=165] (they5.north west);}
\end{tikzpicture}

\hfill{\tiny 0157}\,{\scriptsize \myaudio{./audio/009_answer_be_05.mp3}}

\visible<3->{{\small 答えるときは、名前そのものではなくtheyを用います}}

\end{frame}
%%%%%%%%%%%%%%%%%%%%%%%%%%%
 \subsection{Is this 〜 ?  / Is that 〜 ? と聞かれたら}
%%%%%%%%%%%%%%%%%%%%%%%%%%%
\begin{frame}[plain,t]{Is this  〜 ? / Is that 〜 ? と聞かれたら}
 \large

\begin{columns}[t]
\begin{column}{.49\textwidth}
\visible<2->{Is \myAnch{this}{Maroon}{this} your book?}

\vspace{10pt}

\mbox{}\hfill%
\visible<3->{$\left\{\begin{tabular}{l}
         \visible<3->{Yes, \myAnch{it1}{Maroon}{it} is.}\\
         \visible<4->{No, \myAnch{it2}{Maroon}{it} is not.}\\
         \visible<5->{\text{(}= \text{No, it's not.)}}\\
         \visible<6->{\text{(}= \text{No, it isn't.)}}
       \end{tabular}\right.$}

\begin{tikzpicture}[remember picture,overlay]
 \visible<3->{\draw[->,line width=2pt,opacity=.5,Maroon] (this.south) to[out=-30,in=160] (it1.north west);}
 \visible<4->{\draw[->,line width=2pt,opacity=.5,Maroon] (this.south) to[out=-110, in=165] (it2.north west);}
\end{tikzpicture}
\end{column}
%%%%%%%%%%%%%%
\begin{column}{.49\textwidth}
\visible<7->{Is \myAnch{that}{Maroon}{that} your house?}

\vspace{10pt}

\mbox{}\hfill
\visible<8->{$\left\{\begin{tabular}{l}
         \visible<8->{Yes, \myAnch{it3}{Maroon}{it} is.}\\
         \visible<9->{No, \myAnch{it4}{Maroon}{it} is not.}\\
         \visible<10->{\text{(}= \text{No, it's not.)}}\\
         \visible<11->{\text{(}= \text{No, it isn't.)}}
       \end{tabular}\right.$}

\begin{tikzpicture}[remember picture,overlay]
 \visible<8->{\draw[->,thick,Maroon] (that.south) to[out=-30,in=160] (it3.north west);}
 \visible<9->{\draw[->,thick,Maroon] (that.south) to[out=-90, in=165] (it4.north west);}
\end{tikzpicture}

\hfill{\tiny 0334}\,{\scriptsize \myaudio{./audio/009_answer_be_06.mp3}}

\end{column}
\end{columns}

\visible<3->{{\small 答えるときは、thisやthatではなくitを用います}}

\end{frame}
%%%%%%%%%%%%%%%%%%%%%%%%%%%%%
\begin{frame}<1-20>[plain]\frametitle{Exercises}
例にならって、つぎの質問に対する答えを「はい」と「いいえ」の2通りつくりましょう%
\hfill{\tiny 0109}\,{\scriptsize \myaudio{./audio/009_answer_be_07.mp3}}

\bigskip

\begin{tabular}{rlcll}
\visible<1->{例}& \visible<1->{Are you busy?}& \visible<2->{$\rightarrow$}&\visible<3->{(1) Yes, I am.}&\visible<4->{(2) No, I am not.}\\
\visible<1->{1}&\visible<1->{Are you from Tokyo?}&\visible<5->{$\rightarrow$}&\visible<6->{(1) Yes, I am.}&\visible<7->{(2) No, I am not.}\\
\visible<1->{2}&\visible<1->{Is she a science teacher?}&\visible<8->{$\rightarrow$}& \visible<9->{(1) Yes, she is.}&\visible<10->{(2) No, she is not.}\\
\visible<1->{3}&\visible<1->{Is Peter in Japan now?}&\visible<11->{$\rightarrow$}&\visible<12->{(1) Yes, he is.}&\visible<13->{(2) No, he is not.}\\
\visible<1->{4}&\visible<1->{Is math easy for you?}&\visible<14->{$\rightarrow$}&\visible<15->{(1) Yes, it is.}&\visible<16->{(2) No, it is not.}\\
\visible<1->{5}&\visible<1->{Is that your car?}&\visible<17->{$\rightarrow$}&\visible<18->{(1) Yes, it is.}&\visible<19->{(2) No, it is not.}
\end{tabular}
\end{frame}
%%%%%%%%%%%%%%%%%%%%%%%%%%%%%%%%
\begin{frame}[plain,t]\frametitle{Exercises}
%\begin{ascolorbox3}{Susanの自己紹介を読んで、問に答えましょう}[orange][coltitle=orange!50!black]
\begin{tcolorbox}[colframe=ForestGreen,
  colback=ForestGreen!10!white,
  colbacktitle=ForestGreen!40!white,
  coltitle=black, %fonttitle=\bfseries,
  title=Susanの自己紹介を読んで、問に答えましょう]
\parindent=15pt

\noindent{}Hi, I'm Susan, and I'm from\tikzmark{clue1} New York. Nice to meet you!

Look at this picture. It's my family.
Dad, George, is a \tikzmark{clue2}dentist, and Mom, Helen, tea\tikzmark{clue3}ches.
I have a fifteen-year-old brother, Kevin, and our sweet dog, Brownie.\tikzmark{clue4}

I am really happy to meet you. Thank you!%
\hfill{\tiny 0341}\,{\scriptsize \myaudio{./audio/009_answer_be_08.mp3}}
\end{tcolorbox}
%\end{ascolorbox3}



\begin{tabular}{rll}
1&\visible<2->{Is Susan from New York?}\tikzmark{q1}&\visible<4->{Yes, she is.}\\
2&\visible<2->{Is Susan's father a teacher?}\tikzmark{q2}&\visible<6->{No, he isn't.}\\
3&\visible<2->{Is Helen a dentist?}\tikzmark{q3}&\visible<8->{No, she isn't.}\\
4&\visible<2->{Is the name of Susan's dog Brownie?\tikzmark{q4}\hspace{20pt}\mbox{}}&\visible<10->{Yes, it is.}
\end{tabular}

\mbox{}\hfill\visible<1->{{\scriptsize Nice to meet you. はじめまして\hspace{15pt}look at ～を見る\hspace{15pt}dentist 歯医者\hspace{15pt}Susan's: Susanの所有格}}

\vspace*{-10pt}

\hfill{\tiny 0312}\,{\scriptsize \myaudio{./audio/009_answer_be_09.mp3}}


\begin{tikzpicture}[remember picture,overlay]
 \visible<3->{\draw[<-,opacity=.4,line width=2pt] ([yshift=-2pt]pic cs:clue1) to[bend left] ([xshift=2pt, yshift=2pt] pic cs:q1);}
 \visible<5->{\draw[<-,opacity=.4,line width=2pt] ([yshift=-2pt]pic cs:clue2) to[bend left] ([xshift=2pt, yshift=2pt] pic cs:q2);}
 \visible<7->{\draw[<-,line width=2pt, opacity=.4] ([yshift=-2pt]pic cs:clue3) to[bend left] ([xshift=2pt, yshift=2pt] pic cs:q3);}
 \visible<9->{\draw[<-,opacity=.4,line width=2pt] ([yshift=2pt]pic cs:clue4) to[bend left] ([xshift=2pt, yshift=2pt] pic cs:q4);}
\end{tikzpicture}
\end{frame}
%%%%%%%%%%%%%%%%%%%%%%%%%%%
\begin{frame}[plain]
 
\includegraphics[height=\textheight]{./images/suzans_family.png}
\hfill{\tiny 0341}\,{\scriptsize \myaudio{./audio/009_answer_be_08.mp3}}

\end{frame}
%%%%%%%%%%%%%%%%%%%%%%%%%%
\section{まとめ}
%%%%%%%%%%%%%%%%%%%%%%%%%
\begin{frame}[plain,t]{要点}
\small
\begin{tblr}{
  colspec = {l l l l l},
  row{1} = {font=\bfseries}, % ヘッダーを太字に
}
\toprule
疑問文 & はい (Yes) & いいえ (No) & \shortstack{パターン1 \\ {\scriptsize (S $+$ beを短縮)}} & \shortstack{パターン2 \\ {\scriptsize (be $+$ not を短縮)}} \\
\midrule
{Am I} ...?      & Yes, you are. & No, you are not. & No, \textbf{you're} not.          & No, you \textbf{aren't}.   \\
{Are you} ...?   & Yes, I am.    & No, I am not.    & No, \textbf{I'm} not.      & ---       \\
{Is he} ...?     & Yes, he is.   & No, he is not.   & No, \textbf{he's} not.     & No, he \textbf{isn't}.     \\
Is your father ...?\\
{Is she} ...?    & Yes, she is.  & No, she is not.  & No, \textbf{she's} not.    & No, she \textbf{isn't}.    \\
Is your mother ...?\\
{Is it} ...?     & Yes, it is.   & No, it is not.   & No, \textbf{it's} not.     & No, it \textbf{isn't}.     \\
%{Are we} ...?    & Yes, we are.  & No, we are not.  & No, \textbf{we're not}.    & No, we \textbf{aren't}.    \\
%{Are you} ...? (複数) & Yes, we are. & No, we are not. & No, \textbf{we're} not. & No, we \textbf{aren't}.    \\
{Are they} ...?  & Yes, they are.& No, they are not.  & No, \textbf{they're} not.  & No, they \textbf{aren't}.  \\
Are your parents ...?\\
\bottomrule
\end{tblr}

{\tiny audio\_overview 1740}\,{\scriptsize \myaudio{./audio/overview/009_answer_be_audio_overview.mp4}}

\end{frame}
\end{document}
